\documentclass[9pt, a4paper]{article}

\usepackage[utf8]{inputenc}
\usepackage[T1]{fontenc}
\usepackage[left=10mm, right=10mm, top=10mm, bottom=10mm]{geometry}
\usepackage[default]{lato}
\usepackage{fontawesome5}
\usepackage{hyperref}
\usepackage{xcolor}
\usepackage{titlesec}
\usepackage{enumitem}

\linespread{0.83}

\definecolor{nameblue}{RGB}{31, 73, 125}

\hypersetup{colorlinks=true, urlcolor=black, hidelinks}

\titleformat{\section}{\bfseries\normalsize}{}{0em}{\uppercase}[\vspace{-3pt}\rule{\linewidth}{0.4pt}]
\titlespacing{\section}{0pt}{4pt}{1pt}

\setlength{\parindent}{0pt}
\setlength{\parskip}{0pt}

\newcommand{\workheader}[3]{%
  \textbf{#1}, \textit{#2} \hfill #3%
}

\newcommand{\projheader}[2]{%
  \textbf{#1} \hfill \textit{#2}%
}

\begin{document}
\pagestyle{empty}

%--- HEADER ------------------------------------------------------------------
\begin{center}
  {\Large\bfseries\color{nameblue} Vaibhav Raheja}\\[1pt]
  \small
  \faEnvelope\enspace\href{mailto:vaibhavvraheja@gmail.com}{vaibhavvraheja@gmail.com}
  \quad\faMapMarker*\enspace San Francisco, CA
  \quad\faPhone\enspace +1 (217) 202-9970
  \quad\faLinkedin\enspace\href{https://linkedin.com/in/vaibhav-raheja}{Linkedin}
  \quad\faLink\enspace\href{https://vaibhav-raheja.github.io}{Portfolio}
  \quad\faGithub\enspace\href{https://github.com/Vaibhav-Raheja}{Github}
\end{center}

%--- PROFESSIONAL SUMMARY ----------------------------------------------------
\section{Professional Summary}
Robotics Software Engineer with 2+ years building and deploying autonomous platforms across hardware-software architecture, multi-sensor perception, and production field operations. Strong across ROS2-based autonomy, multi-sensor fusion (LiDAR, thermal, PTZ), and full-stack robotics infrastructure, with field-proven results achieving 90\% uptime and 40\% improvement in fault detection rates.

%--- WORK EXPERIENCE ---------------------------------------------------------
\section{Work Experience}

\workheader{TerraWise Solutions, Inc}{Founding Engineer}{Mar 2026 – Present $|$ San Francisco, CA}
\begin{itemize}[leftmargin=1.2em, topsep=2pt, itemsep=0pt, parsep=0pt]
  \item Building TerraWise's autonomous robotics stack from scratch, establishing ROS2 architecture and Docker deployment workflows following an internal transition from EarthSense.
  \item Defining hardware-software standards and foundational perception pipelines, drawing on multi-sensor integration and field deployment experience from Solarbot.
  \item Migrating the full platform to NVIDIA Jetson, integrating ODrive motor controllers, and implementing LiDAR-based stop-by-obstacle detection.
\end{itemize}

\workheader{EarthSense, Inc}{Robotics Engineer / Project Lead}{Aug 2024 – Mar 2026 $|$ Champaign, USA}
\begin{itemize}[leftmargin=1.2em, topsep=2pt, itemsep=0pt, parsep=0pt]
  \item Led a cross-functional Solarbot team of 8+ engineers to deliver an autonomous solar panel inspection system across 3+ multi-site robot units.
  \item Redesigned the C++ motion controller safety state machine to detect zero-radius turns, enforce a full-stop intermediary between drive modes, and add configurable RPM ramp up/down, fixing misdetection bugs and hardening stop logic for field safety.
  \item Shipped 35+ PRs across the Python/C++ stack spanning motor control, per-motor RPM/current telemetry, camera bringup, and ARM64/Jetson platform deployment, all validated on a HIL rig I designed and built.
  \item Built a low-latency teleoperation system using AWS Kinesis WebRTC (UDP-based video transport) and WebSocket over TCP for command/control, enabling real-time remote monitoring of field robots.
  \item Integrated multi-sensor perception stacks (Hesai/Unitree LiDAR, PTZ cameras, Seek thermal, RealSense) into production via ROS2 and Docker, stress-tested across solar farm deployments in 43--46°C conditions.
\end{itemize}

\workheader{Intelligent Motion Laboratory}{Robotics Research Developer}{Aug 2023 – Dec 2023 $|$ Champaign, USA}
\begin{itemize}[leftmargin=1.2em, topsep=2pt, itemsep=0pt, parsep=0pt]
  \item Achieved sub-2mm tracking accuracy for autonomous robotic eye examinations using MediaPipe FaceMesh and ZED depth cameras on UR5.
  \item Designed a custom camera mount in Fusion 360 for the UR5, improving head tracking precision by 20\%.
\end{itemize}

\workheader{All India Institute of Medical Sciences (AIIMS) Hospital}{Robotics Research Developer, NMIMS}{Feb 2021 – May 2023 $|$ Mumbai, India}
\begin{itemize}[leftmargin=1.2em, topsep=2pt, itemsep=0pt, parsep=0pt]
  \item Engineered a robot-assisted intubation system on xArm5 with HOTAS teleoperation and real-time camera feedback for critical care.
  \item Fabricated a custom catheter end-effector with integrated miniature camera for real-time airway visualization.
\end{itemize}

%--- SKILLS ------------------------------------------------------------------
\section{Skills}
\textbf{Programming:} Python, C++, ROS/ROS2, OpenCV, PyTorch, Machine Learning (ML), CNNs \newline
\textbf{Robotics \& Perception:} LiDAR Integration (Hesai), Multi-Sensor Fusion, Camera Calibration, Path Planning, Motion Planning, Control Algorithms, Reinforcement Learning, Gazebo, SLAM, Docker \newline
\textbf{Networking \& Protocols:} TCP/IP, UDP, WebSocket, WebRTC, CAN, Serial Communication \newline
\textbf{Tools:} Linux, Git, Fusion 360, CAD, Arduino, Raspberry Pi, 3D Printing

%--- PROJECTS ----------------------------------------------------------------
\section{Projects}

\projheader{Benchmarking Control Algorithms for Unitree Go1 Robot}{Python, ISAAC Sim, Reinforcement Learning $|$ 2024}
\begin{itemize}[leftmargin=1.2em, topsep=2pt, itemsep=0pt, parsep=0pt]
  \item Built an ISAAC Sim benchmarking framework for the Unitree Go1, comparing factory vs.\ RL-based locomotion across velocity tracking, stability, and terrain adaptability.
  \item Verified real-world RL controller performance by deploying ``Walk These Ways'' on a physical Go1 across grass, gravel, and inclines.
\end{itemize}

\projheader{GRAIC Autonomous Racing}{Python, Hybrid A*, Path Planning, PD Control, Vehicle Dynamics $|$ 2023}
\begin{itemize}[leftmargin=1.2em, topsep=2pt, itemsep=0pt, parsep=0pt]
  \item Achieved 40.8\% lap time improvement on the Shanghai circuit using a Hybrid A* path planner tuned to vehicle kinematics and non-holonomic dynamics.
  \item Benchmarked three algorithms (Hybrid A*, Dynamic Programming, Spline Interpolation) and tuned a PD controller for steering and velocity.
\end{itemize}

\projheader{Intelligent Ground Vehicle Competition (IGVC), Team DARVIN}{Python, ROS, OpenCV, ODrive, Velodyne LiDAR, Gazebo $|$ 2019--2023}
\begin{itemize}[leftmargin=1.2em, topsep=2pt, itemsep=0pt, parsep=0pt]
  \item Served as Co-Captain, securing 2nd in the Cyber Challenge (2023) and 3rd in the Auto-Nav Challenge (2022) against 15+ international university teams.
  \item Developed SOCRATES 2.0 on ROS with Velodyne LiDAR avoidance, OpenCV lane detection, GPS navigation, and ODrive S1 motor control, Gazebo-validated and field-deployed.
\end{itemize}

%--- EDUCATION ---------------------------------------------------------------
\section{Education}

\textbf{University of Illinois Urbana-Champaign} \hfill Aug 2023 – Dec 2024 $|$ Champaign, USA\\
\textit{Master of Engineering in Autonomy and Robotics}

\vspace{1pt}
\textbf{Mukesh Patel School of Technology Management \& Engineering} \hfill Jul 2019 – Jun 2023 $|$ Mumbai, India\\
\textit{Bachelor of Technology in Computer Engineering}

%--- PUBLICATIONS ------------------------------------------------------------
\section{Publications}
V. Raheja, V. Shah, M. Shetty, and P. Patel, ``Multi-Disease Prediction System using Machine Learning,'' \textit{2022 International Conference on Futuristic Technologies (INCOFT)}, Belgaum, India, 2022, pp. 1--6. doi: 10.1109/INCOFT55651.2022.10094382

\end{document}
